\section{MARCO METODOLÓGICO}
El desarrollo de la presente investigación se apoya en métodos los cuales principalmente están contemplados dentro del método científico, por otra parte se aplican otros tipos de métodos para el diseño y desarrollo del protocolo de votación electrónica que se detallan a continuación:

\subsection{Metodología de Investigación Científica}
El método científico es un proceso dinámico, que envuelve observar todo el tiempo, buscar información continuamente y planificar experimentos para demostrar la hipótesis.

\subsection{Metodología para el Diseño y Desarrollo del Protocolo: Design Science Research}
Para el diseño y desarrollo del protocolo de votación electrónica se emplea \textit{Design Science Research} (DSR), propuesto por Vaishnavi y Kuechler \cite{Vaishnavi2004}, el cual organiza el proceso en cinco fases: Conciencia del Problema, Sugerencia, Desarrollo, Evaluación y Conclusión, admitiendo iteraciones entre ellas cuando los resultados de la verificación formal así lo requieran.

\subsection{Procedimiento General (Fases de Design Science Research)}
La investigación sigue las cinco fases propuestas por Design Science Research, adaptadas al desarrollo y evaluación del protocolo de votación electrónica:


\begin{itemize}
	\item \textbf{Conciencia del Problema.} Se caracteriza la brecha existente entre los estándares FIPS 203 y FIPS 204, publicados en 2024, y su escasa aplicación en protocolos de votación electrónica dentro del contexto universitario boliviano. Esta fase corresponde al planteamiento del problema desarrollado en los Capítulos I y II.
	
	\item \textbf{Sugerencia.} Se formula la propuesta inicial de solución mediante el diseño de un protocolo de votación electrónica basado en criptografía post-cuántica, incorporando los algoritmos ML-KEM y ML-DSA, fundamentados en los problemas matemáticos Module-LWE y Module-SIS, en concordancia con los objetivos establecidos en el Capítulo III.
	
	\item \textbf{Desarrollo.} Se construye el artefacto tecnológico mediante el modelado formal del protocolo en \textit{Tamarin Prover}, considerando el modelo de atacante Dolev--Yao. Asimismo, se formalizan las propiedades de confidencialidad, integridad y autenticidad mediante \textit{lemmas} y se desarrolla una implementación de referencia destinada a la simulación experimental.
	
	\item \textbf{Evaluación.} Se verifica el cumplimiento de las propiedades de seguridad demostrando los \textit{lemmas} definidos previamente y se miden indicadores de rendimiento como latencia, tiempo de ejecución y consumo de memoria bajo diferentes cantidades de votantes simulados. Posteriormente, se comparan estos resultados con un protocolo equivalente basado en RSA/ECDSA mediante un diseño cuasi-experimental. En caso de que alguno de los \textit{lemmas} no se satisfaga o el rendimiento resulte insuficiente, el proceso retorna a la fase de Desarrollo para ajustar el diseño del protocolo.
	
	\item \textbf{Conclusión.} Se consolidan y analizan los resultados obtenidos durante la investigación, los cuales constituyen la base para la prueba de hipótesis, las conclusiones y las recomendaciones finales de la tesis.
\end{itemize}

\subsection{Instrumentos y Herramientas}

El análisis formal del protocolo se realiza mediante \textit{Tamarin Prover}, herramienta ampliamente utilizada para la verificación formal de protocolos criptográficos y aplicada previamente al análisis de sistemas de votación electrónica como Helios, Belenios y Civitas (Baloglu et al., 2021). Su selección responde a su capacidad para razonar sobre protocolos con estado utilizando el modelo de atacante Dolev--Yao.

La implementación de referencia del protocolo se desarrollará utilizando un lenguaje de programación compatible con las bibliotecas criptográficas post-cuánticas seleccionadas, mientras que las métricas de rendimiento se registrarán mediante instrumentos de medición de tiempo de ejecución, consumo de memoria y latencia durante las simulaciones.

La descripción específica del lenguaje de programación, bibliotecas criptográficas, características del hardware utilizado e instrumentos de medición será incorporada una vez definidos estos elementos durante la fase de implementación.
